\chapter{Compléments et repères pour la soutenance}

\section{Paramètres principaux}
\begin{table}[H]
\centering
\caption{Paramètres utilisés dans les notebooks de clustering.}
\label{tab:params}
\begin{tabularx}{\textwidth}{l l X}
\toprule
Méthode & Paramètres & Interprétation \\
\midrule
K-Means & $K=2$, \texttt{init=k-means++}, \texttt{n\_init=20}, \texttt{random\_state=42} & Deux groupes autour de centroïdes. \\
CAH & $K=2$, \texttt{linkage=ward} & Coupe de la hiérarchie donnant la meilleure silhouette testée. \\
DBSCAN & \texttt{eps=0,80}, \texttt{min\_samples=5} & Trois groupes et un label bruit pour 230 clients. \\
\bottomrule
\end{tabularx}
\end{table}

\section{Formules utiles}
\begin{align}
\text{log1p}(x)&=\log(1+x),\\
z&=\frac{x-\mu}{\sigma},\\
IQR&=Q_3-Q_1,\\
s(i)&=\frac{b(i)-a(i)}{\max(a(i),b(i))}.
\end{align}

\section{Lecture rapide des fichiers Python}
\texttt{src/check\_dataset.py} charge le CSV et vérifie dimensions, colonnes, types, valeurs manquantes, doublons et dépenses attendues. \texttt{src/preprocessing.py} centralise l'ordre des six dépenses et construit les pipelines de transformation. Les notebooks utilisent ensuite ces fichiers préparés sans écraser la source.

\section{Glossaire}
\begin{description}
\item[Centroïde] moyenne vectorielle des observations d'un cluster K-Means.
\item[Couverture] proportion des observations effectivement affectées à un cluster ; pour DBSCAN, elle exclut le bruit.
\item[Inertie] somme des distances quadratiques aux centroïdes, critère spécifique à K-Means.
\item[Silhouette] mesure interne combinant proximité au propre cluster et séparation du cluster voisin.
\item[Bruit] observation DBSCAN qui n'est rattachée à aucune zone dense selon les paramètres choisis.
\item[Skewness] mesure d'asymétrie d'une distribution.
\end{description}

\section{Ressources graphiques}
Les 18 figures utilisées dans le rapport sont copiées dans \texttt{reports/latex/figures/}. Les dix tableaux CSV exportés sont conservés dans \texttt{reports/latex/tables/} pour consultation et réutilisation dans Overleaf ou dans une analyse ultérieure.
